\section{Attribution and certificate trust}

Content addressing and attribution answer different questions. A digest can show
that two observers refer to the same bytes, but it does not show who generated
those bytes, under which authority, or whether the declared execution occurred.
Conversely, an authenticated principal can sign a misleading or incomplete
statement. The receipt must bind both content and a scoped attestation.

\begin{definition}[Attribution-complete receipt]
A receipt is attribution-complete for an execution claim only if its attested
payload binds the executor principal, issuer and subject namespace, host/runtime,
provider/model/checkpoint when used, inference or execution configuration,
protocol/prompt digest, tool versions, code revision, environment digest,
evidence-manifest digest, seed, session, campaign, request and result digests.
\end{definition}

This definition is deliberately stronger than requiring a non-empty executor,
model or certificate string. Attestation systems such as in-toto bind supply-chain
steps to authorized actors \citep{torresarias2019intoto}; SLSA provenance records
the build definition and producing platform \citep{slsa2026provenance}; and
identity-based signing can connect a short-lived certificate, an identity and a
transparency log entry \citep{sigstore2026signing}. These are relevant design
donors, not evidence that a research receipt inherits their guarantees.

\begin{definition}[Certificate trust predicate]
For certificate $c$ and policy $T$, $\operatorname{Trust}(c,T)$ holds only when
the signature verifies to a configured root, the issuer and subject namespaces are
allowed, the subject is authorized for the requested role, the validity epoch
covers the execution, revocation and transparency requirements pass, and the
signed payload binds the receipt's canonical bytes.
\end{definition}

An arbitrary identifier fails this predicate even if it is unique. A certificate
for code signing also cannot automatically authorize scientific adjudication; role
and scope are part of the policy. When no appropriate trust root is configured,
the system may preserve the self-declared metadata for audit but must return
\cannotcheck{} for authenticated attribution.

Environment attribution is similarly substantive. Tools such as ReproZip capture
files, libraries and configuration needed to rerun a computational experiment
\citep{chirigati2016reprozip}, while workflow packages can bundle plans, runs,
inputs and outputs \citep{khan2019cwlprov,leo2024workflowrun}. A receipt contract
should interoperate with these representations rather than replace them. Its added
requirement is that the environment record be content-bound to the specific request
and result and assessed against the claim being made.

